\documentclass[11pt]{article}
\usepackage[T1]{fontenc}
\usepackage{lmodern,amsmath,amssymb,amsthm,microtype}
\usepackage[margin=1.07in]{geometry}
\usepackage[hidelinks]{hyperref}
\newtheorem{theorem}{Theorem}
\newtheorem{lemma}[theorem]{Lemma}
\newcommand{\E}{\mathbb E}
\newcommand{\one}{\mathbf 1}
\newcommand{\norm}[1]{\left\|#1\right\|}
\title{A sharp stationary moment threshold away from equality\\for Boolean Doob chains with two bases}
\author{Lluis Eriksson}
\date{19 November 2026}
\begin{document}\maketitle
\begin{abstract}
For Boolean noise of correlation one half, consider a reversible Doob transform whose stationary law factors through a binary linear quotient with two disjoint column bases. Above stationary moment order $15/7$, exact matching of increasing-order coordinate marginals gives full ordered spectral recovery by the previously established packing criterion. For every $2\le p<15/7$ we construct chains in the same class with exact matching through a linear fraction of the coordinates and exponential convergence of the stationary density to one in $L^p$, yet with exponentially small spectral gap. The construction uses ordinary noise of correlation $1/4$ on a duplicated-coordinate quotient. A rare sphere traps the chain while most of the potential is supported near typical magnetizations. A convex correction on the two endpoints of each edge enforces exact moment matching even when the minimum density tends to zero. Conversely, at the second moment a uniformly bounded inverse stationary moment of order above $35/9$ suffices for recovery at arbitrary rank. At equality we prove full ordered spectral recovery for a paired quotient with an arbitrary positive potential depending on a fixed number of block magnetizations, each on a growing block. A layer trace and matching on signed central boxes control all exceptional layers together. The general two-base endpoint remains open.
\end{abstract}

\section{Setting and result}
On $\{-1,1\}^L$ let $\mu$ be uniform and let $P_\rho$ flip coordinates independently with probability $(1-\rho)/2$. For a strictly positive potential $h$ write
\[
 K_h(x,y)=P_\rho(x,y)h(y)/(P_\rho h)(x),\qquad
 \pi_h=hP_\rho h/\E_\mu(hP_\rho h).
\]
The stationary measure is $\nu=\pi_h\mu$. Its reversible edge law is proportional to $\mu(x)P_\rho(x,y)h(x)h(y)$. Matching through degree $k$ means that every nonconstant Walsh character of degree at most $k$ has zero $\nu$-mean. Equivalently, the marginals on any $k$ or fewer coordinates are uniform. Code and matrix terminology uses the identification of signs with bits. The spectral gap is $1-\lambda_2$, with eigenvalues in nonincreasing order and counted with multiplicity.

\begin{theorem}\label{two}
At ambient correlation $1/2$, for each fixed $2\le p<15/7$ there exist $\beta,c>0$ and positive potentials in every sufficiently large even dimension $L=2r$ such that
\[
 \|\pi_h-1\|_{p}\le e^{-cL},\qquad
 \E_\nu\chi_T=0\quad(1\le |T|\le\lfloor\beta L\rfloor),\qquad
 1-\lambda_2(K_h)\le e^{-cL}.
\]
The stationary law factors through the matrix $[I_r\ I_r]$. Its displayed columns form two disjoint bases. In particular the second eigenvalue differs from the ordinary ambient value $1/2$ by at least $1/2-e^{-cL}$.
\end{theorem}

The positive result of \cite[section 7]{companion} applies to any binary quotient admitting two disjoint column bases. For $p>15/7$, a uniformly bounded stationary $L^p$ norm and matching through order $k\to\infty$ imply ordered spectral error $O(e^{-c_p k})$, uniformly in rank. It assumes positivity in each finite dimension. Theorem~\ref{two} proves that its strict moment threshold cannot be lowered throughout this class. The construction allows a vanishing stationary minimum. Theorem~\ref{inverse} below proves recovery at the second moment under a finite inverse moment, in particular under a fixed positive floor.

Our construction works more generally at quotient noise $\rho$. Put
\[
 j=\frac{1-\rho}{1+\rho},\qquad b=\frac{\rho^2}{1-\rho^2}.
\]
\begin{theorem}\label{general}
Fix $0<\rho<1$ and $1<p<\infty$ with $1/p>j-2b$. For all sufficiently large $r$, on the $r$-cube at correlation $\rho$, there are strictly positive potentials for which the three estimates of Theorem~\ref{two} hold with $r$ in place of $L$.
\end{theorem}
At $\rho=1/4$ we have $j=3/5$, $b=1/15$ and $j-2b=7/15$. For the lift to $2r$ coordinates, let $y_i=x_i x_{i+r}$ and pull the quotient potential back along this map. Each $y_i$ experiences independent correlation $(1/2)^2$. The stationary density pulls back exactly. A Walsh coefficient of the pullback vanishes unless its index is a union of complete pairs; complete pairs give exactly the corresponding quotient coefficient. Matching therefore holds to the stated ambient linear order, decreasing $\beta$ if needed. The centered indicator of the inverse image of a quotient trap has the same variance and Dirichlet energy as its quotient indicator. This proves Theorem~\ref{two} from Theorem~\ref{general}.

\section{Ordinary returns and transitions to typical magnetizations}
Write $m(x)=r^{-1}\sum x_i$ and $I(d)=((1+d)\log(1+d)+(1-d)\log(1-d))/2$.
For a small positive constant $d$, let $S$ be the sphere at magnetization $d_r=d+O(1/r)$ obtained by rounding the plus count to the nearest integer. Let $B=\{|m|\le\alpha d\}$ with $0<\alpha<\rho$, and write $q=P_\rho(x,S)$ for $x\in S$, independent of $x$ by symmetry. We use three estimates:
\begin{align}
 \mu(S)&\le e^{-rI(d)+O(1)},\quad
 q\ge \exp[-rJ(d)-O(\log r)], \label{eq:return}\\
 \sup_{S} P_\rho(x,B)&\le
 \exp\left[-r\left(\frac{(\rho-\alpha)^2}{1-\rho^2}\frac{d^2}{2}
                  -O_\rho(d^3)\right)+O(1)\right],\label{eq:bulk}\\
 \sup_{B} P_\rho(y,S)&\le \exp[-r(1-\rho\alpha)^2d^2/2+O(1)] .\label{eq:back}
\end{align}
Here $J(d)=(1+\rho)I(2d/(1+\rho))/2-I(d)=j I(d)+O_\rho(d^4)$. Bounds are applied with $d,\alpha$ fixed before the dimension grows.

For \eqref{eq:return}, the first estimate is the usual binomial bound. In a stationary noise edge the probabilities of $++,+-,-+,--$ are $(a,z,z,a)$ with $a=(1+\rho)/4$ and $z=(1-\rho)/4$. Choose the same nearest integer to $rz$ for each off-diagonal count and fix the diagonals by the two prescribed sphere margins. Stirling's formula bounds the probability of this type below by $\exp[-r(1+\rho)I(2d/(1+\rho))/2-O(\log r)]$. All four limiting cell proportions are positive. Divide by $\mu(S)$ and use transitivity. This return estimate also occurs in the earlier shell construction \cite{shell}. For \eqref{eq:bulk}, condition on the input $x$. The output signs have means $\rho x_i$ and each variance $1-\rho^2$. Their centered log moment generating functions at $-u$ are $(1-\rho^2)u^2/2+O_\rho(u^3)$ uniformly for small $u$. Apply the exponential Markov bound with $u=(\rho-\alpha)d/(1-\rho^2)$. Finally, from $B$ the mean output magnetization is at most $\rho\alpha d$. Hoeffding's bound for independent signs at threshold $d_r$ gives \eqref{eq:back}.

\section{A trap with negligible stationary moment cost}
Choose a constant $\ell$ satisfying
\[
 \max(0,(j-1/p)/2)<\ell<\min(j,b).
\]
Such a choice follows from the strict hypothesis and $j>0$. Choose $\alpha$ small enough that $(\rho-\alpha)^2/(1-\rho^2)>\ell$ and $(1-\rho\alpha)^2>j-\ell$. Since $I(d)=d^2/2+O(d^4)$, choose $d>0$ small enough that the following inequalities all have strict exponential slack:
\[
 J(d)-2\ell I(d)<I(d)/p,\quad
 \sup_S P_\rho(x,B)\le e^{-r(\ell I(d)+\eta)},\quad
 e^{r(J(d)-\ell I(d))}\sup_B P_\rho(y,S)\le e^{-\eta r}
\]
for some $\eta>0$ and all large $r$. The set $B$ has complement of exponentially small uniform measure.

Put $P=P_\rho$, $H_0=1$ on $B$ and $H_0=e^{-Dr}$ elsewhere, where $D$ is a fixed constant larger than $2+|J(d)|+\ell I(d)$. Set
\[
 M=e^{-\ell I(d)r}/q,\qquad H=H_0+M\one_S,\qquad W=H P H.
\]
For $x\in S$, $P H_0(x)/(Mq)\le e^{-c r}$ for some $c>0$. The numerator bounds all contributions from outside $S$. Consequently each row of the transformed kernel starting on $S$ has exit probability at most $e^{-c r}$.

We verify $\norm{W-1}_p\le C e^{-c r}$, allowing $c$ to decrease from one estimate to the next. The background obeys $\norm{H_0 P H_0-1}_p\le2\mu(B^c)^{1/p}$, by contraction of $P$ and $0<H_0\le1$. Off $S$ the extra term is $H_0 M P\one_S$. On $B$ it is exponentially small by \eqref{eq:return} and \eqref{eq:back}, absorbing the polynomial factor in the slack. Elsewhere off $S$ it is at most $e^{-Dr}M$, also exponentially small. On $S$ we have $PH\le 2Mq$ and $H\le2M$ eventually (use $q\le1$). Hence
\[
 \norm{\one_S W}_p
 \le 4\,\mu(S)^{1/p}\frac{e^{-2\ell I(d)r}}{q}
 \le \exp\{r[J(d)-2\ell I(d)-I(d)/p]+O(\log r)\}.
\]
The stationary density $W/\E W$ therefore converges to one exponentially in $L^p$. The very small background outside $B$ makes an additive correction to the density potentially incompatible with positivity. We correct the potential itself.

\section{Exact marginals by tilting both ends of an edge}
\begin{lemma}\label{tilt}
Let $P$ be a reversible positive semidefinite Markov operator for a finite uniform measure $\mu$, let $H>0$, and let $\pi$ be its normalized stationary density as above. Let $\phi_1,\ldots,\phi_N$ be distinct nonconstant orthonormal Walsh characters on a cube. If $N\varepsilon$ is sufficiently small, where $\varepsilon=\norm{\pi-1}_1$, there exists $f$ in their span with $\norm f_\infty\le C N\varepsilon$ such that $He^f$ has zero stationary mean for every listed character. Constants are absolute.
\end{lemma}
\begin{proof}
Use the normalized symmetric edge law $J$ of $H$. For $\Phi=(\phi_i)_i$ define $\psi(s)=\log \E_J e^{s\cdot(\Phi(X)+\Phi(Y))}$. Its gradient at zero has Euclidean norm at most $2\sqrt N\,\varepsilon$. The tilt by $s$ is exactly the edge law of $H e^{s\cdot\Phi}$, with marginal $\nu_s$. For any Euclidean unit $v$, put $F=v\cdot\Phi$. Then $|F|\le\sqrt N$, $\E_\mu F=0$, $\E_\mu F^2=1$, and
\[
 v^\top\nabla^2\psi(s)v
 =2\operatorname{Var}_{\nu_s}F+
    2\operatorname{Cov}_{J_s}(F(X),F(Y))
 \ \ge\ 2\operatorname{Var}_{\nu_s}F.
\]
The covariance is nonnegative: the centered quadratic form of the edge law is that of $P$ applied to the potential times the centered function, divided by its normalization. At zero, the variance is at least $1-N\varepsilon-N\varepsilon^2$. In the Euclidean ball of radius $1/(4\sqrt N)$ the marginal density relative to $\nu_0$ is at least $e^{-1}$, by the two endpoint factors and their normalization. The variance comparison follows by minimizing the squared distance to a constant. Thus the Hessian is bounded below by an absolute positive multiple of the identity throughout this ball. For $N\varepsilon$ sufficiently small, Taylor's inequality gives $\psi(s)>\psi(0)$ on its boundary. The minimum is interior, and strong convexity bounds its distance from zero by $C\sqrt N\,\varepsilon$. At that minimum the gradient is twice the vector of stationary character means. Taking $f=s\cdot\Phi$ proves the bound.
\end{proof}
Use all characters of degrees $1$ through $\lfloor\beta r\rfloor$. Their number is at most $\exp[r H_{\rm bin}(\beta)]$ for $0<\beta<1/2$ (natural entropy). Choose $\beta$ sufficiently small using the exponential $L^1$ estimate. Then $\|f\|_\infty$ is exponentially small. The new stationary density differs from the old by a pointwise factor between $e^{-4\|f\|_\infty}$ and $e^{4\|f\|_\infty}$. The $L^p$ conclusion follows. Each transition changes by at most $e^{2\|f\|_\infty}$ as a ratio; the small exit probability survives. The corrected stationary mass of $S$ is at most $\|\pi_{\rm new}\|_p\mu(S)^{1-1/p}=o(1)$. Testing the gap variational formula with the centered indicator of $S$ gives a gap at most $\sup_{x\in S}K(x,S^c)/(1-\nu_{\rm new}(S))$. This proves Theorem~\ref{general}.

\section{Recovery with a finite inverse stationary moment}
The counterexample uses a deep depression of the stationary density. A finite inverse moment suffices to exclude it even when the pointwise minimum tends to zero. We record the positive result in the same two-base class.

\begin{theorem}\label{inverse}
Let the ambient correlation be $1/2$. Suppose that $\pi$ factors through a linear quotient with two disjoint full column bases, matches through degree $k$, and satisfies $\|\pi\|_2\le B$ and $\|\pi^{-1}\|_s\le D$ for fixed $s>35/9$. Then the maximum difference between corresponding ordered eigenvalues of $K_h$ and $P_{1/2}$ is at most $C e^{-ck}$, with constants depending only on $B,D,s$. In particular a fixed positive floor suffices at the second moment for any rank.
\end{theorem}
\begin{proof}
Normalize $h$ by $\E hP h=1$, and set $P=P_{1/2}$, $Q=P_{1/\sqrt2}$, $g=Ph$, $F=Qh$, $U=h/g=\pi/g^2$. Then $\|F\|_2=1$. The positive scaling $h$ inherits the quotient symmetry from $\pi$ by uniqueness: the positive diagonal of $P$ makes the stationary-scaling functional strictly convex in $\log h$ \cite{shell}. Two disjoint bases give the quotient contraction $P_\lambda:L^a\to L^{1+(a-1)\lambda^{-4}}$ for $a>1$. This is ordinary forward hypercontractivity in the first basis followed by the same estimate in the second, with any remaining column noise a contraction \cite{companion}. Thus $\|g\|_5=\|QF\|_5\le1$. Jensen gives $g^{-1}=1/P(\pi/g)\le P(g/\pi)$. Set $1/a=1/5+1/s$. Then
\[
 \|g^{-1}\|_{q}\le D,\qquad q=1+16(a-1)>20.
\]
This gain is the reason for the inverse moment hypothesis. By H\"older, $U$ is bounded in $L^v$ for $1/v=1/2+2/q<3/5$.

For clarity, the comparison from moments to the spectrum can be carried out on the full cube. Put $a_0=\E h$ and $V=\|F-a_0\|_2^2=1-a_0^2$. Cauchy--Schwarz gives $1=(\E hg)^2\le a_0\E(\pi g)$. Since $g=QF$ and $\pi-1$ has only degrees $>k$, we have $V\le a_0 B\,2^{-(k+1)/2}\sqrt V$. Thus $\|g-1\|_2\le C_B2^{-k/2}$. On $E=\{|g-1|>2^{-k/8}\}$, uniform measure is at most $C_B2^{-3k/4}$. Consequently $\|(U-\pi)\one_E\|_{5/3}\le C\mu(E)^{3/5-1/v}$ (here $v<2$). Off $E$, $|U-\pi|\le C2^{-k/8}\pi$ for large $k$.

For each nonnegative quotient weight $W$, $\|QM_W Q\|\le\|W\|_{5/3}$ on the full cube. Indeed, its positive leading eigenvector for $W\ne0$ is invariant under kernel translations by uniqueness, so the quotient $L^2\to L^5$ estimate bounds its quadratic form. It follows that $\|QM_{U-\pi}Q\|$ decays exponentially. For the signed comparison choose $\lambda=3/4$, $\theta=1/(\sqrt2\lambda)<1$. Since $1+\lambda^{-4}\ge4$, the same positive multiplier argument bounds $\|P_\lambda M_{\pi-1}P_\lambda\|$ by $\|\pi-1\|_2$ using $|\pi-1|$. Split each $P_\theta$ in $Q=P_\theta P_\lambda$ at degree $\lfloor k/2\rfloor$. The low--low form vanishes by matching, and each other block carries a decaying factor $\theta^{\lfloor k/2\rfloor+1}$. Thus $\|QM_UQ-P\|$ is exponentially small. The reversible conjugate of $K_h$ is $M_{\sqrt U}P M_{\sqrt U}$. It has the same eigenvalues as $QM_UQ$ by the two Gram products of $M_{\sqrt U}Q$. Min--max completes the comparison.
\end{proof}
The number $35/9$ is sufficient; we do not assert necessity. In particular the traps of Theorem~\ref{two} at $p=2$ cannot have a uniformly bounded inverse stationary moment of any fixed order above $35/9$.

\section{The critical exponent for block-symmetric paired coordinates}
\label{sec:critical}
On the paired quotient $y_i=x_ix_{r+i}$ the ambient $P=P_{1/2}$ acts as $P_\rho$ with $\rho=1/4$. Split the $r$ pairs into $m$ blocks of sizes $n_i$. Write $d_i$ for the mean sign in block $i$ and $z_i=\sqrt{n_i}d_i$. A block-symmetric potential is any positive function of the vector $(d_1,\ldots,d_m)$; the function may couple the blocks arbitrarily.

\begin{theorem}\label{critical}
Fix $m$ and let $\min_i n_i\to\infty$. If the potential is block-symmetric, $\|\pi\|_{15/7}\le B$ uniformly, and the stationary ambient marginals match through $k\to\infty$, then
\[
 \max_j|\lambda_j(K_h)-\lambda_j(P_{1/2})|\longrightarrow0 .
\]
There is no lower-tail or uniform-integrability hypothesis on the stationary density. For $m=1$ the subcritical examples of Theorem~\ref{two} can all be chosen radial.
\end{theorem}
For the last assertion the untilted potential and the set of constraints in Lemma~\ref{tilt} are invariant under coordinate permutations. The strictly convex local minimization there has a unique minimizer in its invariant ball. Its tilt is therefore invariant too.

\paragraph{Two rates for one block.}
For a block of $n$ signs let $\mu_d$ denote the uniform mass of the layer $d$ and $T(d,s)$ the ordinary quotient transition between layers. Let $I_2(d,s)$ be the Legendre transform of the joint log moment-generating function of one uniform stationary noise pair. Define
\[
 J(d)=I_2(d,d)-I(d),\qquad E(d)=I_2(d,0)-I(d).
\]
These are the conditional rates at $d$ and zero for a row starting at $d$. Indeed, for dual parameter $t$ on the output, optimizing over the input parameter gives the binary marginal cost $I(d)$ and leaves the conditional log moment-generating function
\[
 L_d(t)=\tfrac{1+d}{2}\log(\cosh t+\rho\sinh t)
       +\tfrac{1-d}{2}\log(\cosh t-\rho\sinh t).
\]
This follows by tilting the input Bernoulli marginal to mean $d$, with endpoint values by continuity. In particular $E(d)=-\min_t L_d(t)$.
At $p=15/7$ set $\Delta(d)=J(d)-2E(d)-I(d)/p$. Then
\begin{equation}\label{delta}
 \Delta(d)\ge0,\qquad \inf_{|d|\ge1/2}\Delta(d)>0.
\end{equation}
To check this, write $v=(1+\rho)/2$, $j=(1-\rho)/(1+\rho)$ and $b=\rho^2/(1-\rho^2)$. The diagonal rate is the transform of $\log(1-v+v\cosh t)$: by symmetry the two dual parameters can be taken equal. This function is at most $v^{-1}\log\cosh(vt)$, with strictly increasing difference (larger minus smaller) for $t>0$. Differentiation reduces the latter assertion to $(1-v)\sinh(vt)>v\sinh((1-v)t)$. Inserting the optimizer $t=v^{-1}\operatorname{arctanh}d$ of the comparison transform shows that $J\ge jI$, with a uniform additive gain for $|d|\ge1/2$, including the endpoints by continuity.

For the other direction apply the two-function reverse hypercontractive inequality for binary noise \cite{mos} to indicators of layers at $d$ and zero. For $a,b'\in(0,1)$ with $(1-a)(1-b')=\rho^2$, their joint probability is at least $\mu_d^{1/a}\mu_0^{1/b'}$. On even block lengths approximate $d$ by admissible layer values. The joint exponential rate is $I_2(d,0)$: enumerating the four-cell types gives the minimum relative entropy with these margins, since there are polynomially many types and a minimizing table can be rounded with the prescribed margins and $O(1/n)$ entry errors. Stirling's formula gives the marginal rates $I(d)$ and zero. Taking the dimension limit first, then $b'\downarrow0$, gives $E\le bI$. Both rate functions are continuous on the closed interval by the same finite table formula. Now $j-2b=1/p$ proves \eqref{delta}.

\paragraph{Layer transition estimates.}
Uniformly for the lattice values with $|d|\le1/2$,
\[
 \mu_d\ge c n^{-1/2}e^{-nI(d)},\qquad T(d,d)\le C n^{-1/2}e^{-nJ(d)} .
\]
The first estimate is Stirling's formula. For the second, tilt the independent output signs to mean sum $nd$. The tilt is bounded because the target stays in a compact subinterval of $(-1,1)$ and both initial success probabilities are bounded away from zero and one. At least one of the two sign groups has size $\ge n/2$. Its tilted binomial count has maximum atom $O(n^{-1/2})$, and convolution preserves this upper bound. The change of measure on the target layer is exactly $e^{-nJ(d)}$. For arbitrary $d$ the weaker bounds $\mu_d\ge(n+1)^{-1}e^{-nI(d)}$ and $T(d,d)\le e^{-nJ(d)}$ follow from the binomial type bound and Chernoff's inequality, including end layers.

For $1\le \sqrt n\,|s|\le2$ on the same side of zero as $d$ (either side at zero), we also have
\begin{equation}\label{lowerrow}
 T(d,s)/\mu_s\ \ge\ c\exp[-nE(d)+\rho\sqrt n\,|d|]
\end{equation}
uniformly in $d$. Tilt this time to mean zero. The optimizing parameter $t_0$ has opposite sign to $d$ and $|t_0|\ge\rho|d|$, since $L'_d(0)=\rho d$ and $L''_d\le1$. Also $|t_0|\le\operatorname{arctanh}\rho$ by the extreme all-plus and all-minus inputs. Under this tilt the probability of each specified layer in the window is at least $c/\sqrt n$. Indeed the smaller sign group lies within $2\sqrt n$ of its tilted mean with probability at least $3/4$ by Chebyshev. Conditional on such a count, the larger group needs a count within $C\sqrt n$ of its mean and of the correct lattice parity. Stirling's bounds give a uniform lower atom $c/\sqrt n$ there, since its size is at least $n/2$ and its tilted parameter stays in a compact subinterval of $(0,1)$. Undo the tilt, multiplying by $\exp[-nE(d)-nt_0s]$, and use $\mu_s\le C/\sqrt n$.

\paragraph{The central mass supplied by matching.}
Normalize $\E hPh=1$, and keep $g=Ph$, $Q=P_{1/\sqrt2}$, $U=\pi/g^2$. Under the stationary law the vector $z$ converges in distribution to $m$ independent standard Gaussians. For each fixed total degree, every polynomial moment is eventually the uniform one because a quotient coordinate is a product of two ambient signs. The independent uniform block sums have Gaussian limiting joint moments; tightness, uniform higher moments and Gaussian moment determinacy give the assertion. In particular each of the $2^m$ signed boxes $1\le |z_i|\le2$ has stationary probability bounded below. Each box also has uniformly positive integral of $h$: $\E g=\E h\le1$ by $\|Qh\|_2=1$, and $\E[\pi\one_{g>L}]\le B L^{-1/p'}$ for $p'=p/(p-1)$. Choose $L$ fixed so this is smaller than half the box mass and use $h=\pi/g$. Multiplying \eqref{lowerrow} for the signed box matching the row $d$ and integrating gives
\begin{equation}\label{gbox}
 g(d)\ge c_{m,B}\exp\!\left[-\sum_i n_i E(d_i)+\rho\sum_i|z_i|\right].
\end{equation}
In particular $g$ is bounded below on each fixed central box $\max_i|z_i|\le R$ because $E(d)=O(d^2)$ at zero (Taylor expansion of $L_d$ at its bounded optimizer near zero).

\paragraph{A tail trace on the full ambient space.}
For any nonnegative orbit weight $W$, the top eigenvector of $QM_WQ$ can be taken invariant under the fiber translations and the within-block permutations. For $W\ne0$ the matrix has strictly positive entries and its Perron eigenvector is unique after normalization. Its norm is thus bounded by the trace on the space of functions of the block magnetizations. On this space $Q^2=P$ has the product transition on layers. Cyclicity of trace, or the orthonormal basis of normalized layer indicators, gives
\[
 \|QM_WQ\|\le\sum_d W(d)\prod_i T(d_i,d_i).
\]
For $W=U\one_{\max |z_i|>R}$ apply H\"older against $\prod_i\mu_{d_i}$, using \eqref{gbox}. The $p'$-power of the resulting multiplier norm in H\"older is bounded by a constant times the tail sum of the product of
\[
 A_n(d)=\mu_d^{1-p'}T(d,d)^{p'}
       \exp[2p'n E(d)-2p'\rho\sqrt n\,|d|].
\]
For $|d|\le1/2$ the two local prefactors give
$A_n(d)\le C n^{-1/2}e^{-p'n\Delta(d)-2p'\rho|z|}$.
For $|d|>1/2$ the weaker estimates and the strict gap in \eqref{delta} give $A_n(d)\le(n+1)^{p'-1}e^{-c n}$. Summing over the mesh $2/\sqrt n$ in $z$ gives a uniformly bounded one-block sum with tail $O(e^{-cR})+o_{n\to\infty}(1)$. The union bound on the product, with fixed $m$, proves
\[
 \limsup_{\min n_i\to\infty}\|QM_{U\one_{\max|z_i|>R}}Q\|
 \le C_{m,B}e^{-cR}.
\]

\paragraph{Completing the spectral comparison.}
On the fixed central box \eqref{gbox} bounds $g^{-2}$ and the matching variance estimate in Theorem~\ref{inverse}'s proof gives $g\to1$ in uniform probability. Thus $\|(U-\pi)\one_{\max|z_i|\le R}\|_{5/3}\to0$: apply H\"older to $\pi(g^{-2}-1)$ with $p>5/3$ and use bounded convergence in probability for the second factor on the box. On the complementary tail $\|\pi\one_{\max|z_i|>R}\|_{5/3}\le B\,\mu(\max |z_i|>R)^{3/5-1/p}$, which tends uniformly to zero as $R\to\infty$ by the uniform-law sign tail bound. The positive multiplier bound from Theorem~\ref{inverse} applies to both terms; the trace estimate treats the remaining $U$ tail. First taking the dimension limit and then $R\to\infty$ yields $\|QM_{U-\pi}Q\|\to0$. The signed comparison $\|QM_{\pi-1}Q\|\to0$ from the same proof uses only matching and bounded second moment. The two Gram products and min--max give Theorem~\ref{critical}. General two-base quotients and growing numbers of blocks are outside this endpoint argument.


\section{Scope and antecedents}
The rare-sphere joint type and calibration by its return probability are shared with \cite{shell}. There, compensation bands preserve a uniform floor and the all-correlation conclusion reaches some moment above $(1+\rho)/(1-\rho)$. Here the potential on the bulk complement can be arbitrarily small. This permits the full one-step cost $b$ and the range $1/p>j-2b$. The correction in Lemma~\ref{tilt} works with the edge law even when the density has a very small minimum. On a quotient obtained by repeating each coordinate $t$ times, ambient correlation $1/2$ gives $\rho=2^{-t}$ and $j-2b=(2^t-1)/(2^t+1)-2/(4^t-1)$. For unrestricted two-base quotients our sharp-threshold claim concerns fixed $p\ge2$ away from equality. Theorem~\ref{critical} includes equality for paired quotients with finitely many permutation blocks. At equality the proof has no strict interval for $\ell$. The available rank-dependent recovery criterion \cite{enumerator} is consistent with these traps; it imposes an additional condition on the rank relative to the matching order.

Classical results of Kirshner and Samorodnitsky \cite{ks} describe near extremality of spheres under ordinary Boolean noise. Reverse hypercontractivity \cite{mos} underlies the earlier sufficient packing estimate. These are inputs and comparisons for finite chains. Our result concerns the transformed stationary law, exact marginals and the gap of its actual kernel. We make no literature-wide priority claim.

This version extends \href{https://airr.science/papers/ARR-2026-54GRS2WX6182EVEB/versions/v1/}{the first version} with Theorem~\ref{critical}. Classical ordinary-noise estimates do not cover the nonlinear denominator of the transformed kernel.

\section*{Model assessments and review context}
OpenAI Codex assisted with the construction, proof and manuscript in the programme research chat. Any assessment accompanying a deposit refers to the bytes and version identified in that record. The arguments here are analytic. The author administers AIRR.

\begin{thebibliography}{9}
\bibitem{companion} L. Eriksson, \emph{Stationary concentration and local return in Boolean Doob chains}, \href{https://airr.science/papers/ARR-2026-2SH3CGGBHG9Q5BYN/versions/v1/}{AIRR working paper, v1}, section 7.
\bibitem{shell} L. Eriksson, \emph{A metastable Hamming shell with exactly matched stationary marginals}, \href{https://airr.science/papers/ARR-2026-2B2W1WQ7DF9N5VW7/versions/v1/}{public v1}; \href{https://github.com/lluiseriksson/THE-ERIKSSON-PROGRAMME/blob/1b24d36de2f0596e915fc605a86d52dc25bdf8c8/papers/stationary-weighting354/research/2026-11-18-small-support/metastable-shell-v2.tex}{subsequent source, section 6}.
\bibitem{enumerator} L. Eriksson, \emph{Controlling the denominator by the weight enumerator}, \href{https://github.com/lluiseriksson/THE-ERIKSSON-PROGRAMME/blob/0fa67e0bb0636fad345c58da03d9426c05d4deca/papers/stationary-weighting354/research/2026-11-18-small-support/enumerator-denominator.tex}{source revision}, 2026.
\bibitem{ks} N. Kirshner and A. Samorodnitsky, \emph{A moment ratio bound for polynomials and some extremal properties of Krawchouk polynomials and Hamming spheres}, \href{https://arxiv.org/abs/1909.11929}{arXiv:1909.11929}.
\bibitem{mos} E. Mossel, K. Oleszkiewicz and A. Sen, \emph{On reverse hypercontractivity}, \href{https://arxiv.org/abs/1108.1210}{arXiv:1108.1210}.
\end{thebibliography}
\end{document}
